\documentclass[11pt]{article}
\usepackage[margin=1in]{geometry}
\usepackage{amsmath,amssymb,amsthm}
\usepackage[hidelinks]{hyperref}

\newtheorem*{theorem}{Literature Answer}

\title{Negative Literature Answer to the Bounded Multiplier Question in arXiv:1206.3358}
\author{FA/Banach run packet}
\date{2026-07-18}

\begin{document}
\maketitle

\section*{Status}

This is a literature-already-answered packet. It records that the bounded
Fourier multiplier equality question from Chen--Xu--Yin \cite{CXY} was answered
negatively in the later paper of Ricard \cite{Ricard}. It is not an original
proof from this run.

\section*{Original Problem}

In Section 7, after proving the theta-independence of completely bounded
\(L_p\) Fourier multipliers, Chen--Xu--Yin ask:

\begin{quote}
Let \(2<p\leq\infty\). Does one have
\[
\mathrm{M}(L_p(\mathbb T^d_\theta))=\mathrm{M}(L_p(\mathbb T^d))?
\]
\end{quote}

They conjecture a negative answer and give only an infinite-generator
\(p=\infty\) illustration.

\section*{Supporting Answer}

Ricard \cite{Ricard} explicitly starts from the Chen--Xu--Yin completely
bounded result and states that bounded multipliers behave differently. In the
arXiv source, Theorem \texttt{ex1} proves that for any irrational
\(\theta\) and any \(1<p\neq 2<\infty\) there is
\(\phi:\mathbb Z^2\to\mathbb C\) such that
\[
\|M_\phi:L_p(\mathbb T^2)\to L_p(\mathbb T^2)\|<\infty
\quad\text{but}\quad
\|M_\phi:L_p(\mathbb T^2_\theta)\to L_p(\mathbb T^2_\theta)\|=\infty.
\]
This disproves the equality for every \(2<p<\infty\).

For the endpoint \(p=\infty\), Proposition \texttt{ex3} in \cite{Ricard}
constructs, for every irrational \(\theta\), a symbol \(\phi:\mathbb Z^2\to
\mathbb C\) such that
\[
\|M_\phi:L_\infty(\mathbb T^2_\theta)\to L_\infty(\mathbb T^2_\theta)\|<\infty
\quad\text{but}\quad
\|M_\phi:L_\infty(\mathbb T^2_\theta)\to L_\infty(\mathbb T^2_\theta)\|_{cb}
=\infty.
\]
Ricard recalls immediately before this proposition that the latter cb norm
agrees with the classical \(L_\infty(\mathbb T^2)\) multiplier norm by the
Chen--Xu--Yin theorem. Hence the endpoint equality is also false.

\begin{theorem}
The equality
\[
\mathrm{M}(L_p(\mathbb T^d_\theta))=\mathrm{M}(L_p(\mathbb T^d))
\]
asked in \cite{CXY} has a negative answer in the finite-dimensional quantum
torus setting for every \(2<p\leq\infty\). It already fails for irrational
two-tori.
\end{theorem}

\begin{proof}[Literature identification]
For \(2<p<\infty\), apply Ricard's Theorem \texttt{ex1} with \(d=2\) and
irrational \(\theta\). The displayed symbol is a bounded classical multiplier
but is not bounded on the corresponding quantum \(L_p\)-space, so the two
multiplier spaces cannot be equal.

For \(p=\infty\), apply Ricard's Proposition \texttt{ex3}. The constructed
symbol is bounded on \(L_\infty(\mathbb T^2_\theta)\) but has infinite
completely bounded norm there. Since the completely bounded norm is equal to
the classical \(L_\infty(\mathbb T^2)\) multiplier norm, the same symbol is not
a bounded classical multiplier. Thus the endpoint equality also fails.
\end{proof}

\section*{Scope Limitations}

This packet records a known negative answer for irrational two-dimensional
quantum tori. It does not classify rational theta or the exact structure of
the bounded multiplier spaces.

\section*{Verification Notes}

The source and supporting papers are distinct. The supporting paper is a later
paper whose abstract and main results explicitly distinguish bounded
multipliers from the completely bounded theta-independent result of
Chen--Xu--Yin. The local duplicate search covered arXiv ids
\texttt{1206.3358} and \texttt{1512.01142}, plus phrases
\texttt{bounded Fourier multipliers}, \texttt{quantum tori},
\texttt{M(L\_p(T\_theta))}, and \texttt{Ricard}.

\begin{thebibliography}{9}

\bibitem{CXY}
Z. Chen, Q. Xu, and Z. Yin,
\emph{Harmonic analysis on quantum tori},
arXiv:1206.3358; Commun. Math. Phys. 322 (2013), 755--805.

\bibitem{Ricard}
E. Ricard,
\emph{\(L_p\)-Multipliers on Quantum Tori},
arXiv:1512.01142; J. Funct. Anal. 270 (2016), 4604--4613.

\end{thebibliography}

\end{document}
